\documentclass[10pt,a4paper]{amsart}
\usepackage{amsmath,amssymb,amsthm,graphicx}
% Retained preamble for the Stacks Project chapter {CH} (shared project preamble).
% The chapter file pinned-original-source.tex only carries \input{{preamble}} on its first line,
% so no preamble text lives inside the chapter itself.  The Stacks Project uses one single shared
% preamble for every chapter: the snapshot stages it three times (sources/stacks-chow/preamble.tex,
% sources/stacks-etale-cohomology/preamble.tex and focused-seeds/AGP/sources/stacks-snapshot-20261011/
% preamble.tex) and those copies are identical in size.  The block below is copied mechanically,
% line for line, from sources/stacks-chow/preamble.tex lines 14-18 and 154-217: the Xy-pic setup and
% the native theorem declarations together with the mathematical macros.  Nothing was added, removed
% or redefined; the wrapper class and the amsmath/amssymb/amsthm/graphicx base packages are supplied
% by the assembler.
% For commutative diagrams we use Xy-pic
% Chapter: stacks-spaces-more-morphisms.tex -- More on Morphisms of Spaces.
\usepackage[all]{xy}

% We use 2cell for 2-commutative diagrams.
\xyoption{2cell}

% Theorem environments.
%
\theoremstyle{plain}
\newtheorem{theorem}[subsection]{Theorem}
\newtheorem{proposition}[subsection]{Proposition}
\newtheorem{lemma}[subsection]{Lemma}

\theoremstyle{definition}
\newtheorem{definition}[subsection]{Definition}
\newtheorem{example}[subsection]{Example}
\newtheorem{exercise}[subsection]{Exercise}
\newtheorem{situation}[subsection]{Situation}

\theoremstyle{remark}
\newtheorem{remark}[subsection]{Remark}
\newtheorem{remarks}[subsection]{Remarks}

\numberwithin{equation}{subsection}

% Macros
%
\def\lim{\mathop{\mathrm{lim}}\nolimits}
\def\colim{\mathop{\mathrm{colim}}\nolimits}
\def\Spec{\mathop{\mathrm{Spec}}}
\def\Hom{\mathop{\mathrm{Hom}}\nolimits}
\def\Ext{\mathop{\mathrm{Ext}}\nolimits}
\def\SheafHom{\mathop{\mathcal{H}\!\mathit{om}}\nolimits}
\def\SheafExt{\mathop{\mathcal{E}\!\mathit{xt}}\nolimits}
\def\Sch{\mathit{Sch}}
\def\Mor{\mathop{\mathrm{Mor}}\nolimits}
\def\Ob{\mathop{\mathrm{Ob}}\nolimits}
\def\Sh{\mathop{\mathit{Sh}}\nolimits}
\def\NL{\mathop{N\!L}\nolimits}
\def\CH{\mathop{\mathrm{CH}}\nolimits}
\def\proetale{{pro\text{-}\acute{e}tale}}
\def\etale{{\acute{e}tale}}
\def\QCoh{\mathit{QCoh}}
\def\Ker{\mathop{\mathrm{Ker}}}
\def\Im{\mathop{\mathrm{Im}}}
\def\Coker{\mathop{\mathrm{Coker}}}
\def\Coim{\mathop{\mathrm{Coim}}}

% Boxtimes
%
\DeclareMathSymbol{\boxtimes}{\mathbin}{AMSa}{"02}

%
% Macros for moduli stacks/spaces
%
\def\QCohstack{\mathcal{QC}\!\mathit{oh}}
\def\Cohstack{\mathcal{C}\!\mathit{oh}}
\def\Spacesstack{\mathcal{S}\!\mathit{paces}}
\def\Quotfunctor{\mathrm{Quot}}
\def\Hilbfunctor{\mathrm{Hilb}}
\def\Curvesstack{\mathcal{C}\!\mathit{urves}}
\def\Polarizedstack{\mathcal{P}\!\mathit{olarized}}
\def\Complexesstack{\mathcal{C}\!\mathit{omplexes}}
% \Pic is the operator that assigns to X its picard group, usage \Pic(X)
% \Picardstack_{X/B} denotes the Picard stack of X over B
% \Picardfunctor_{X/B} denotes the Picard functor of X over B
\def\Pic{\mathop{\mathrm{Pic}}\nolimits}
\def\Picardstack{\mathcal{P}\!\mathit{ic}}
\def\Picardfunctor{\mathrm{Pic}}
\def\Deformationcategory{\mathcal{D}\!\mathit{ef}}
\usepackage{cleveref}
\crefname{originallemma}{original lemma}{original lemmas}
\Crefname{originallemma}{Original lemma}{Original lemmas}
\makeatletter
\newlabel{sites-modules-lemma-localize-differentials}{{04BO}{source}{}{sourcepointer-1}{}}
\newlabel{sites-modules-lemma-localize-differentials@cref}{{[originallemma][0][]04BO}{[1][0][]source}{}{}{}}
\makeatother
\crefname{originallemma}{original lemma}{original lemmas}
\Crefname{originallemma}{Original lemma}{Original lemmas}
\makeatletter
\newlabel{spaces-properties-equation-restrict}{{03LQ}{source}{}{sourcepointer-2}{}}
\newlabel{spaces-properties-equation-restrict@cref}{{[originallemma][0][]03LQ}{[1][0][]source}{}{}{}}
\makeatother
\crefname{originallemma}{original lemma}{original lemmas}
\Crefname{originallemma}{Original lemma}{Original lemmas}
\makeatletter
\newlabel{spaces-properties-lemma-etale-morphism-topoi}{{03LP}{source}{}{sourcepointer-3}{}}
\newlabel{spaces-properties-lemma-etale-morphism-topoi@cref}{{[originallemma][0][]03LP}{[1][0][]source}{}{}{}}
\makeatother
\title{The sheaf of differentials of a morphism of algebraic spaces restricts to the relative differentials along an etale commutative square}
\author{The Stacks Project Authors}
\date{stacks-project:spaces-more-morphisms.tex:local-snapshot-20261011; GNU FDL 1.2 or later}
\begin{document}
\maketitle
\noindent\textit{Editorial note.} All delivered mathematical text is the author's own, copied from the pinned Stacks Project chapter file spaces-more-morphisms.tex; no statement or proof step was rewritten, reordered or removed.\par\medskip
\noindent\textit{Editorial note.} Preamble provenance, as required for this source: the first line of the pinned chapter file is \textbackslash{}input\{preamble\}, and the native preamble was not staged with the source, so no preamble text lives inside the chapter. The retained preamble used here is the project's single Stacks-wide preamble (the same 217-line / 8116-byte file that the already published chow and etale-cohomology cards use), copied mechanically line for line from sources/stacks-chow/preamble.tex lines 14-18 (Xy-pic setup) and 154-217 (theorem declarations and mathematical macros). Not one character was changed, no package was added, no macro was defined, and the excerpt serves only as preamble, never as a span coordinate.\par\medskip
\noindent\textit{Editorial note.} The same-chapter pointers are located at their \textbackslash{}label lines inside the pinned chapter file; the cross-chapter pointers are located in the project's own tag list staged in the snapshot (sources/stacks-chow/tags.txt), at the verbatim lines carrying their canonical tags.\par\medskip
\noindent\textit{Allowed standard prerequisites (editorial).}\begin{itemize}
\item Etale morphisms of algebraic spaces induce equivalences of small etale topoi. Properties of Spaces, Lemma labelled spaces-properties-lemma-etale-morphism-topoi and Equation labelled spaces-properties-equation-restrict; invoked at lines 999-1003. Turns restriction of a sheaf on X\_etale to U\_etale into pullback along a\_\{small\}, and likewise for b.
\item Localisation of the module of differentials on ringed sites. Modules on Sites, Lemma labelled sites-modules-lemma-localize-differentials; invoked at lines 1003-1005. Computes Omega\_\{X/Y\}|\_\{U\_etale\} as the module of differentials of the map of sheaves of rings induced by a\_\{small\} and f\_\{small\}.
\end{itemize}\medskip
\ifdefined\hypertarget\hypertarget{sourcepointer-1}{}\fi
\noindent\textit{Source reference.} Original lemma 04BO: sources/stacks-chow/tags.txt, source lines 5659--5659. Navigation pointer for \textbackslash{}ref\{sites-modules-lemma-localize-differentials\} used inside the counted span (line 1004); the argument that consumes it is retained in the span.\par\medskip
\ifdefined\hypertarget\hypertarget{sourcepointer-2}{}\fi
\noindent\textit{Source reference.} Original lemma 03LQ: sources/stacks-chow/tags.txt, source lines 4722--4722. Navigation pointer for \textbackslash{}ref\{spaces-properties-equation-restrict\} used inside the counted span (line 1001); the argument that consumes it is retained in the span.\par\medskip
\ifdefined\hypertarget\hypertarget{sourcepointer-3}{}\fi
\noindent\textit{Source reference.} Original lemma 03LP: sources/stacks-chow/tags.txt, source lines 4721--4721. Navigation pointer for \textbackslash{}ref\{spaces-properties-lemma-etale-morphism-topoi\} used inside the counted span (line 1000); the argument that consumes it is retained in the span.\par\medskip
\begin{lemma}
\label{lemma-localize-differentials}
Let $S$ be a scheme. Let $f : X \to Y$ be a morphism of algebraic spaces
over $S$. Consider any commutative diagram
$$
\xymatrix{
U \ar[d]_a \ar[r]_\psi & V \ar[d]^b \\
X \ar[r]^f & Y
}
$$
where the vertical arrows are \'etale morphisms of algebraic spaces. Then
$$
\Omega_{X/Y}|_{U_\etale} = \Omega_{U/V}
$$
In particular, if $U$, $V$ are schemes, then this is equal to the usual
sheaf of differentials of the morphism of schemes $U \to V$.
\end{lemma}

\begin{proof}
By
Properties of Spaces, Lemma \ref{spaces-properties-lemma-etale-morphism-topoi}
and Equation (\ref{spaces-properties-equation-restrict})
we may think of the restriction of a sheaf on $X_\etale$ to
$U_\etale$ as the pullback by $a_{small}$. Similarly for $b$. By
Modules on Sites, Lemma \ref{sites-modules-lemma-localize-differentials}
we have
$$
\Omega_{X/Y}|_{U_\etale} =
\Omega_{\mathcal{O}_{U_\etale}/
a_{small}^{-1}f_{small}^{-1}\mathcal{O}_{Y_\etale}}
$$
Since $a_{small}^{-1}f_{small}^{-1}\mathcal{O}_{Y_\etale}
= \psi_{small}^{-1}b_{small}^{-1}\mathcal{O}_{Y_\etale}
= \psi_{small}^{-1}\mathcal{O}_{V_\etale}$ we see that the lemma holds.
\end{proof}

\end{document}
